\chapter{Introduction}
\label{cap:introduction}

\chapterquote{Many things we can't do are simply because we think we can't do them}{Satoru Iwata}

% The study of large musical corpora constitutes one of the primary research tools within disciplines such as historical musicology, ethnomusicology, musical analysis, or sound heritage documentation. Through the systematic analysis of extensive sets of works, it is possible to identify stylistic patterns, recurring compositional structures, relationships between repertoires, and cultural phenomena that can hardly be observed through the isolated study of individual pieces.

% Traditionally, this type of research has depended on manual processes of score reading, cataloging, and analysis—a rigorous methodology but limited in terms of scalability. The growing volume of digitized musical material, especially in fields related to traditional music and musical heritage, has highlighted the need to develop tools capable of automating part of these tasks and facilitating the exploration of large documentary collections.

% The emergence of standardized symbolic formats for musical representation, such as MusicXML or MEI, has driven the development of computational musicology and has enabled the application of automated analysis techniques over increasingly broad repertoires. However, practice demonstrates that significant challenges persist regarding corpus heterogeneity, format compatibility, the extraction of complex musicological information, the integration between musical and textual content, and efficient knowledge retrieval from large volumes of data.

% In parallel, recent advances in artificial intelligence and natural language processing have opened new possibilities for interacting with extensive collections of data. Nevertheless, the particularities of musical information pose limitations that hinder the direct application of these technologies, especially when queries require knowledge derived from specialized musical analyses that are not explicitly represented in the original sources.

% In this context, this project proposes the design and development of an architecture oriented towards the automated analysis of large and heterogeneous musical corpora. The system combines score processing techniques, musical feature extraction, automatic cataloging, information retrieval, and assistance through language models, with the aim of providing a platform capable of facilitating musicological research and improving access to the knowledge contained within large musical collections.

% Within this framework, the project \textit{"\textnormal{\textit{EA}}-DIGIFOLK-An European and Ibero-American approach for the digital collection, analysis and dissemination of folk music (101086338)”} is developed. Funded by the European Union, with the participation of the UCM\footnote{\url{https://cordis.europa.eu/project/id/101086338}}, its primary objective is to preserve and share European and Ibero-American popular folk culture, especially within the musical domain. EA-DIGIFOLK seeks to create a platform that facilitates the dissemination of this culture to educational environments and audiences. Furthermore, and more closely related to this project, another core objective is the musical and ethnomusicological analysis of the collected audiovisual and written material. This work has been carried out as part of this research project.

The study of large musical corpora constitutes one of the main research tools within disciplines such as historical musicology, ethnomusicology, musical analysis, and the documentation of sound heritage. Through the systematic analysis of extensive collections of works, it is possible to identify stylistic patterns, recurrent compositional structures, relationships between repertoires, and cultural phenomena that are difficult to observe through the isolated study of individual pieces.

Traditionally, this type of research has relied on manual processes of reading, cataloguing, and analysing scores, a rigorous methodology but limited in terms of scalability. Likewise, this analogue approach (and even some conventional structured digital approaches) introduces a significant degree of methodological rigidity; when new research questions or unforeseen hypotheses arise, it becomes necessary to re-catalogue the entire corpus and to restructure the relevant database architecture in order to incorporate new variables. The growing volume of digitised musical material, particularly in areas related to traditional music and musical heritage, has highlighted the need to develop tools capable of automating part of these tasks, increasing the flexibility of information usage, and facilitating the dynamic exploration of large documentary collections.

The emergence of standardised symbolic formats for musical representation, such as \textit{MusicXML} \citep{good2001musicxml} or \textit{MEI} \citep{roland2002music}, has fostered the development of computational musicology and enabled the application of automated analysis techniques to increasingly large repertoires. However, practice shows that significant challenges remain, including the heterogeneity of corpora, format interoperability, the extraction of complex musicological information, the integration of musical and textual content, and the efficient retrieval of knowledge from large volumes of data.

In parallel, the emergence of large language models (LLMs) and recent advances in natural language processing have opened a new paradigm in interacting with extensive data collections. Unlike traditional indexed systems, LLMs offer the potential to soften the methodological rigidity of conventional databases, allowing researchers to formulate queries in natural language without the need to manually restructure relational schemas or exhaustively recatalogue the corpus from scratch. However, the particularities of symbolic musical information pose severe limitations that hinder the direct application of these models. Due to the density of formats such as MusicXML \citep{good2001musicxml} or MEI \citep{roland2002music}, the direct injection of scores rapidly saturates model context windows. Moreover, lacking native capabilities for structured computation, LLMs are ineffective when addressing queries that require specialised musical analysis (such as the detection of syncopations, hemiolas, or musical motifs) that are not explicitly represented in the original textual sources, which tends to generate factual hallucinations and restricts their utility in technical musicological research unless they are integrated into hybrid and specialised architectures.

In this context, this work presents the design, development, and implementation of a hybrid AI-based engineering architecture aimed at the automated analysis and flexible querying of large and heterogeneous folk music corpora. To overcome the bias of pure generative technologies, the system deploys an ingestion pipeline capable of processing symbolic collections (\textit{MusicXML}, \textit{MEI}, and \textit{MSCZ}), extracting and normalising, using the \textit{music21} library, a spectrum of high-complexity variables ranging from basic metadata to phenomena requiring more advanced musical analysis, such as syncopations, horizontal and vertical hemiolas, musical motifs, and structural quality control. The originality of the proposed solution lies in the design and implementation of a local, deterministic, hallucination-free infrastructure based on a relational \textit{SQLite} database, interconnected with a small language model \textit{(Llama3.1:8B)} through the advanced \textit{Model Context Protocol} (MCP) standard \citep{anthropic2024mcp} for the automatic translation of natural language into structured queries (\textit{Text-to-SQL}). To evaluate performance and validate the effectiveness of this approach, the research includes a rigorous comparative study against a RAG-based implementation using the native \textit{File Search} engine of \textit{Gemini}, empirically demonstrating how the MCP-structured local approach overcomes the semantic limitations and hallucination rates of large-scale commercial models in complex musicological queries. This solution is exposed through an accessible web application interface, providing the ethnomusicological community with a tool for dynamic exploration and information extraction in large musical collections.

To address the challenges of analytical flexibility, automation of symbolic format processing, and the limitations of generative AI, this architecture is framed within and developed as part of the European and Ibero-American project \textit{“EA-DIGIFOLK — An European and Ibero-American approach for the digital collection, analysis and dissemination of folk music (101086338)”}, funded by the European Union and with active participation from the Complutense University of Madrid (UCM)\footnote{\url{https://cordis.europa.eu/project/id/101086338}}. While the purpose of \textit{EA-DIGIFOLK} is to preserve, share, and spread folk musical heritage in public and educational environments through an accessible platform, its fundamental methodological core lies in the deep musicological and ethnomusicological analysis of the collected audiovisual and written material. It is in this specialised scientific dimension that this research work finds its purpose, providing the technological infrastructure necessary to transform the digitised corpus into explicit, auditable, and dynamic knowledge.

\section{Motivation}
% The musicological analysis of large volumes of traditional and/or folk music has historically faced a methodological limitation: its dependence on the manual analysis of scores, which represents a rigorous but hardly scalable process, making the digital automation of this procedure virtually impossible. While the digitization of musical archives into structured symbolic formats (such as MusicXML or MEI) has opened the doors to computational musicology, fundamental problems that hinder any automation methodology when digitizing this type of content continue to persist.

% The following are some obstacles and difficulties encountered in the field of musical analysis that have motivated the development of this project:

The motivation driving this research lies in the need to address the disconnection between the time and effort required of musicological researchers and the unmanageable scale of today’s digital heritage. Behind each folk music corpus there are ethnomusicologists and sound archivists who devote years to the preservation and understanding of this cultural heritage. However, a significant portion of their time is ultimately absorbed by repetitive tasks of transcription, annotation, and manual cataloguing. Although necessary, these activities limit the time and energy that could otherwise be devoted to analysis, critical interpretation, and the generation of new knowledge about the musical traditions under study.

Traditional ethnomusicological analysis requires exhaustive and highly technical musical examination. Identifying rhythmic structures such as the hemiola, tracing a recurring \textit{leitmotif}, or auditing the quality of a transcription in a folk songbook containing thousands of pieces requires the human eye to traverse countless pages of scores, measure by measure. This analogue approach is rigorous but fundamentally lacks scalability. When the volume of digitised material reaches thousands of works (as is the case in international preservation projects such as \textit{``EA-DIGIFOLK''}), manual analysis ceases to be a slow methodological option and instead becomes a physical and intellectual barrier. In this reality, many musicologists are forced to work with reduced samples that represent only a small fraction of the available repertoire. This limitation not only constrains the questions that can be posed but also reduces the statistical robustness of their conclusions. As a consequence, a large portion of the potential knowledge contained in these archives remains unexplored.

To this physical limitation is added a persistent cognitive frustration derived from the rigidity of current data storage systems. Scientific research is not a linear process but a dynamic one, in which the discovery of new patterns and details generates new questions. Within the conventional paradigm, if a musicologist formulates a hypothesis midway through their research that requires correlating vocal tessitura with a specific lyrical theme, they are confronted with a technological barrier. Evaluating this new variable entails reopening every score in the corpus one by one to extract the missing information, or alternatively relying on a technical development team to restructure the relational database and implement new analysis methods. This technical dependency disrupts the flow of scientific thinking and constrains the researcher’s creativity within the limits of pre-existing database schemas.

Therefore, the motivation of this work is not to replace the expert’s judgement with automated computation, but rather to free them from operational burden. There is an urgent need to design tools that communicate through natural language with humanists, acting as bridges that enable the study of complex collections without requiring code writing or database redesign. By delegating the algorithmic extraction of variables and the translation of complex queries to a deterministic artificial intelligence architecture, musicologists are returned their most valuable resource: time to think, interpret, and safeguard the intangible cultural richness of our communities.

% \begin{itemize}
% \item \textbf{Heterogeneity in the corpora.} Large-scale musical corpora, such as the one used throughout this project, tend to be highly heterogeneous both structurally and in their cataloging. First, works may be stored in multiple file formats (\textit{MSCZ}, \textit{MusicXML}, \textit{MEI}, among others), each with technical and semantic particularities that require specific extraction and normalization processes. This diversity prevents the construction of fully automated analysis workflows and forces the development of preprocessing procedures adapted to each information source. Furthermore, corpus heterogeneity is not limited to file formats but also affects the cataloging criteria and organization of repositories. In these cases, collections are frequently encountered where certain fundamental metadata, such as the title of the work, the author, or the geographical origin, are stored within the structure of the file itself, whereas in other cases, such information is only available in the filename or in the hierarchical organization of the directories. The lack of homogeneous standards for documentary description complicates the automatic identification of works and requires additional mechanisms for inference and metadata validation.

% \item \textbf{Inconsistency in symbolic standards.} Many standard tools used in the technical field of musicological analysis, such as \textit{music21} can experience critical failures or lose relevant metadata when interpreting certain information sources. This occurs, for example, with certain tokens present in MEI files (such as the \textit{wordpos} attribute, which indicates the position of a syllable within a word; however, not all the values it can take are interpretable by many programs like \textit{music21}), whose implementation is not fully normalized among different score editing software.

% \item \textbf{Algorithmic complexity in the detection of complex rhythmic phenomena.} In the world of music, there is an infinity of expressive and ambiguous structures, such as syncopations and hemiolas (both horizontal and vertical), which require mathematical modeling with a strict preliminary study based on very specific characteristics of the works, such as attack times and metric hierarchies. This entails working from scratch on functions that are not natively implemented in conventional analytical libraries.

% \item \textbf{Division between musical and semantic analysis.} Many automated musical analysis systems isolate structural and rhythmic content from the textual content of the lyrics. This is a highly problematic separation, as decoupling these two deeply related structures results in a severe loss of information. Therefore, it is necessary to design a system capable of extracting analytical musical characteristics while simultaneously classifying the thematic and poetic content of the songs.

% \item \textbf{Contextual limitations of language models in extensive corpora.} When the size of the corpus is large, providing the model with all the information as input context becomes unfeasible, particularly in smaller models whose context limits and reasoning capabilities are more restricted. Even when employing information retrieval strategies (\textit{Retrieval-Augmented Generation}, RAG), many musicological queries require the prior extraction and calculation of complex analytical features from the scores, such as rhythmic patterns, melodic structures, or metric relationships. Since these processes cannot be reliably executed in real time by a language model, an intermediate architecture capable of performing specialized musical analysis and storing its results in optimized structures for subsequent retrieval and querying becomes necessary.
% \end{itemize}

% The project originates from these gaps, aiming to resolve them through the study and development of a system capable of meeting the needs of musicologists and users who require more sophisticated tools to handle the analysis of very extensive and complex corpora of data. The main purpose of this research centers on developing an automated, robust, and scalable architecture capable of analyzing latent features in a corpus of musical works, simultaneously integrating structural, rhythmic, melodic, and semantic information to facilitate the musicological study of large folk collections.

\section{Objectives}
% Based on the limitations identified in current computational musicological analysis methodologies, this project seeks to develop a technological infrastructure that facilitates the study of large musical corpora through automated analysis techniques, information retrieval, technical musicological analysis, and natural language processing. To this end, the following objectives are established.

Based on the critical limitations identified in current musicological analysis methodologies (such as the structural rigidity of conventional databases, the ineffectiveness of general-purpose language models when confronted with the density of symbolic formats, and the lack of scalability in manual analysis), this project aims to design and deploy a scalable and accessible technological infrastructure for researchers. This hybrid architecture seeks to unify deterministic musical analysis techniques with the flexibility of conversational artificial intelligence.

\subsection{General Objective}
% To develop an automated, robust, and scalable system for the analysis, cataloging, and querying of large, heterogeneous musical corpora, capable of extracting musicological knowledge from a large number of works stored in different formats, thereby facilitating research, exploration, and information retrieval for musicologists, researchers, and other specialized users.

The general objective of this work is to design, develop, and implement a hybrid, locally deployed, and robust AI-based engineering architecture aimed at automating analysis, systematic cataloguing, and dynamic information retrieval in large corpora of folk and popular music. This system is intended to process and normalise large collections of digital scores in symbolic formats in order to extract various complex musicological variables; these technical and structural descriptors are subsequently organised and modelled within a centralised relational database, which acts as the reliable source that renders the system deterministic. The ultimate goal of this infrastructure is to democratise access for non-expert users to this type of architecture and to facilitate the workflow of musicologists and researchers, freeing them from mechanical cataloguing tasks and traditional programming barriers through a natural language conversational interface. Furthermore, in order to validate the feasibility and technical accuracy of this proposal against the state of the art, the project aims to empirically evaluate and compare the performance of this local structured approach against the capabilities of a large-scale commercial cloud-based model, demonstrating how the combination of relational databases and specialised context protocols can neutralise factual hallucinations and memory saturation issues that large generative systems suffer when performing specialised musical analysis.

\subsection{Specific Objectives}
% \begin{itemize}
% \item Design an architecture capable of processing musical corpora composed of different symbolic representation formats, guaranteeing compatibility between diverse documentary sources.
% \item Develop an automatic system for the extraction and normalization of musical metadata in order to facilitate the systematic cataloging of works and the construction of reliable and consistent databases.
% \item Implement musical analysis algorithms that allow for the detection and quantification of structural, melodic, rhythmic, and metric characteristics relevant to musicological research.
% \item Integrate musical analysis with the semantic and textual analysis of the lyrics, enabling the joint study of both the musical and literary aspects of the works.
% \item Automate processing and analysis tasks that usually require a high investment of time or specialized technical knowledge, reducing barriers to access advanced musical research tools.
% \item Design and implement an information retrieval system that allows for complex searches over the corpus, including queries based on both metadata and automatically extracted musical features.
% \item Evaluate the capabilities and limitations of general-purpose language models in tasks involving the analysis and querying of extensive musicological corpora, identifying the scalability, precision, and context problems that these systems present.
% \item Propose complementary mechanisms to overcome the limitations observed in these models through the incorporation of specialized musicological analysis and optimized structures for knowledge retrieval.
% \end{itemize}

To ensure the execution of the main purpose of this research, the following specific objectives are outlined:

\begin{itemize}
\item \textbf{Multi-format Processing and Ingestion.} In line with the overall objective of supporting heterogeneous musical corpora, the project defines the need to design and implement a modular ingestion pipeline capable of parsing, validating, and unifying symbolic collections encoded in the \textit{MusicXML}, \textit{MEI}, and \textit{MSCZ} (\textit{MuseScore}) standards. This subsystem must manage encoding discrepancies across different documentary sources while ensuring a homogeneous internal representation for subsequent analytical stages.
\item \textbf{Extraction and Normalisation of Structural Metadata.} Derived from the general objective of automating cataloguing processes, this objective involves the development of algorithmic methods for systematically extracting, normalising, and structuring relevant metadata and musical information from each work (such as key signature, meter, tessitura, geographical provenance, and related attributes). The goal is to eliminate inconsistencies associated with manual records and to build a robust relational database.
\item \textbf{Technical Musicological Modelling and Parameterisation.} In support of the broader objective of conducting comprehensive musicological analysis and extracting specialised knowledge, this objective focuses on the implementation of advanced algorithms (leveraging the analytical and music-processing capabilities of the \textit{music21} library) dedicated to detecting, measuring, and quantifying complex abstract musical variables that remain inaccessible to generic AI systems. These include phenomena such as syncopation patterns, hemiolas, tessitura, recurrent melodic profiles, and related structural characteristics.
\item \textbf{Hybrid Musical-Textual Analysis.} Following the general objective of enabling a comprehensive exploration of musical heritage, this objective combines the analysis of musical information contained within scores with the semantic and thematic study of the lyrics associated with folk songs. This integration will allow researchers to relate specific musical characteristics (such as melodic, rhythmic, or harmonic structures) to the themes, narratives, and meanings embedded in oral traditions, facilitating richer and more contextualised interpretations of the repertoire.
\item \textbf{Optimised Relational Persistence Design.} Derived from the objective of ensuring a deterministic and scalable environment, this objective involves the design of a relational database architecture based on \textit{SQLite}. The database engine must be structured with appropriate indexing strategies and integrity constraints capable of efficiently storing the full spectrum of extracted analytical variables, thereby serving as the system’s single source of truth.
\item \textbf{Implementation of the Local Communication Protocol (MCP).} In support of the objective of enabling flexible querying without requiring programming expertise, this objective consists of deploying and configuring a server based on the \textit{Model Context Protocol} (MCP) standard. This protocol will act as a secure and standardised communication layer connecting the relational database with the local language model, enabling the system to interact dynamically with its underlying data environment.
\item \textbf{Orchestration of a Local Text-to-SQL Conversational Engine.} Derived from the objective of enabling natural-language interaction, this objective focuses on the design of model instructions (\textit{prompt engineering}) and the configuration of a locally deployed small language model (\textit{Llama 3.1:8B}). Its purpose is to specialise the model in the real-time translation of complex musicological questions formulated in natural language into precise SQL queries, thereby ensuring accurate responses while minimising the risk of hallucinations.
\item \textbf{Comparative Study and Empirical Evaluation Against Commercial Systems.} In line with the objective of critically assessing the capabilities of language models, this objective consists of conducting a rigorous empirical study comparing the performance of the proposed local structured architecture against a commercial cloud-based baseline using RAG through \textit{Gemini File Search}. The evaluation focuses on measuring system accuracy when handling highly specialised and technically demanding musicological queries.
\item \textbf{Development of a Web Interface and User Experience Layer.} Derived from the objective of democratising access to advanced computational tools, this objective involves the development of a custom web application that encapsulates the technical complexity of the system, its databases, and its MCP-based connections behind an intuitive graphical user interface, eliminating technical barriers for end users.
\end{itemize}

\section{Work Plan}
The work plan to be followed to achieve the objectives described in the previous section is detailed here. In order to attain the proposed objectives, a work plan structured into six main phases has been designed, spanning from late October to early June. This schedule encompasses everything from the initial research and requirements gathering through field work, to the technical development of the architecture, its comparative evaluation, and the final writing of the research report. Below, the temporal planning is detailed using a Gantt chart (Figure~\ref{fig:gantt_plan}):

\begin{figure}[H]
\raggedright
\hspace*{-2.5cm}
\begin{ganttchart}[
    vgrid,
    hgrid,
    x unit=0.05cm,
    y unit chart=0.7cm,
    time slot format=isodate-yearmonth,
    title/.style={fill=gray!20, draw=black, font=\small\bfseries},
    bar/.style={fill=blue!40, draw=black},
    milestone/.style={fill=red, draw=black},
    progress label text={},
    group right shift=0,
    group top shift=.6,
    group height=.2
]{2025-10}{2026-06}
    \gantttitlecalendar{year, month=letter} \\
    
    \ganttgroup{F1: State of the Art \& Initial Res.}{2025-10}{2025-12} \\
    \ganttbar{Literature review}{2025-10}{2025-11} \\
    \ganttbar{Format analysis (MusicXML, MEI)}{2025-11}{2025-12} \\
    
    \ganttgroup{F3: Architectural Design}{2025-12}{2026-02} \\
    \ganttbar{SQLite DB schema design}{2025-12}{2026-01} \\
    \ganttbar{Normalization pipeline development}{2026-01}{2026-02} \\
    
    \ganttgroup{F4: Implementation \& Analysis}{2026-02}{2026-04} \\
    \ganttbar{Musical analysis modules (music21)}{2026-02}{2026-03} \\
    \ganttbar{Semantic analysis integration (Ollama)}{2026-03}{2026-04} \\
    
    \ganttgroup{F2: Requirements Elicitation}{2026-03}{2026-03} \\
    \ganttbar{Interviews with musicologists}{2026-03}{2026-03} \\
    \ganttbar{Definition of user requirements}{2026-03}{2026-03} \\
    
    \ganttgroup{F5: Evaluation \& Testing}{2026-04}{2026-05} \\
    \ganttbar{Tests with Gemini 2.5 Flash (RAG)}{2026-04}{2026-05} \\
    \ganttbar{Comparative evaluation Llama3 vs Gemini}{2026-05}{2026-05} \\
    
    \ganttgroup{F6: Final Documentation}{2026-04}{2026-06} \\
    \ganttbar{Technical report writing}{2026-04}{2026-06} \\
    \ganttmilestone{Report Submission}{2026-06}
\end{ganttchart}
\caption{Gantt chart of the work plan (October - June).}
\label{fig:gantt_plan}
\end{figure}

At an operational level, the development is broken down into the following key activities directly linked to the project's objectives:
\begin{itemize}
    \item \textbf{Phase 1 (October -- December):} Dedicated to the deep study of the state of the art in computational musicology.
    \item \textbf{Phase 2 (December -- February):} Focused on the design and implementation of the intermediate infrastructure and the automatic normalization pipeline, responsible for unifying data persistence into the local SQLite file regardless of the original heterogeneity or organization of the folk corpus.
    \item \textbf{Phase 3 (February -- April):} Focused on the development of algorithms based on \textit{music21} for extracting complex technical features (rhythmic and melodic structures) and their integration with local language models (\textit{Llama3} via \textit{Ollama}) intended for the semantic and thematic classification of the works' lyrics.
    \item \textbf{Phase 4 (March):} Exclusively dedicated to requirements gathering through qualitative interviews with musicologists and researchers in the field. The goal of this milestone is to align the system's capabilities with the actual daily analytical needs of expert users.
    \item \textbf{Phase 5 (April -- May):} Allocated to the critical evaluation phase. Tests are conducted on the information retrieval system, implementing both the local text-to-SQL approach and the cloud-based RAG approach via the native \textit{File Search} functionality of the Gemini API, enabling a contrast of precision, context, and hallucination boundaries between both technologies.
    \item \textbf{Phase 6 (April -- June):} Consisting of writing the research report and preparing the documentation for the developed software to ensure full technical reproducibility later on.
\end{itemize}

\section{Methodology and Technology}
To develop this project, open-source tools and the free version of the Gemini API have been selected. This decision ensures the full reproducibility of the project locally (with the exception of utilizing the Gemini API).
\begin{itemize}
\item \textbf{Python 3.} Used as the base to develop the entire project. This high-level, general-purpose programming language was selected primarily because it hosts the most advanced ecosystem of libraries for computational musicology, data science, and artificial intelligence. Its versatility allows for articulating within a single automated workflow everything from score preprocessing and metadata normalization from heterogeneous musical corpora, to local database management, communication with language model servers via structured requests, and the subsequent analytical evaluation of results.

\item \textbf{Music21.} A leading toolkit for computational musicology. It allows for a deep musical analysis of each work regarding its notes, intervals, chords, scales, etc., which facilitates the study of harmonic, melodic, and rhythmic patterns that are difficult to identify with other tools and, in this case, with LLMs. Additionally, it supports multiple formats, such as MIDI or XML, which is highly useful in this research, where we work with an extensive and heterogeneous corpus of works in different formats. Moreover, its integration with Python and other libraries facilitates the extraction and analysis of information from the musical corpus.

\item \textbf{Ollama.} In this project, Ollama is deployed as a local server responsible for orchestrating the small language model \textit{Llama 3.1 (8B)}. This infrastructure was selected, first, because it guarantees full control over data and absolute privacy for heritage corpora by processing information locally, and second, because it eliminates the dependency on token-based costs and network latency associated with commercial cloud services. Within the proposed architecture, Ollama integrates directly with the advanced \textit{Model Context Protocol} (MCP) and accesses a previously constructed \textit{SQLite} relational database containing information extracted from the musical works of the corpus. The local LLM does not read the scores directly; instead, it operates as a deterministic algorithmic translator (\textit{Text-to-SQL}). It receives natural-language questions posed by researchers and, through carefully designed instructions and structured typing, generates SQL queries that are executed against the database.

\item \textbf{Gemini API (Free Version).} A cloud-based generative artificial intelligence platform developed by Google. It represents the technological counterpart used to evaluate the validity of the proposed local architecture. In this project, the commercial model \textit{Gemini 2.5 Flash} is accessed through its official API (under the development access plan) in order to implement an alternative RAG-based system using its native \textit{File Search} engine. Under this approach, the musical corpus files are uploaded directly to Google's semantic index in the cloud. The methodological role of the Gemini API is strictly comparative: it serves as the \textit{baseline} or industry-standard commercial solution against which the proposed system is evaluated. This comparison aims to demonstrate empirically how a large-scale model with billions of parameters and extensive context windows may still suffer from factual hallucinations, latent operational costs, and contextual misunderstandings when confronted with highly specialized ethnomusicological queries, in contrast to the precision achieved by a small local model coordinated through MCP.

\item \textbf{SQLite.} SQLite constitutes the relational engine and deterministic \textit{ground truth} upon which the local query infrastructure is built. Unlike traditional client-server database management systems such as \textit{MySQL}, \textit{PostgreSQL}, or \textit{SQL Server}, which introduce significant technical and operational overhead through deployment, network configuration, port management, and access-control requirements, \textit{SQLite} operates as an embedded relational database engine. This means that the database library is integrated directly into the application code, consolidating the entire schema, indexes, and musicological variables extracted from the corpus into a single local binary file. This choice offers several advantages. First, it provides optimal read performance for the \textit{Text-to-SQL} workflow managed through MCP, enabling immediate access without network latency. Second, it guarantees complete portability of the ecosystem. By storing all information within a single self-contained file, SQLite removes installation barriers for ethnomusicology research teams, allowing the analysed corpus and its query engine to be transferred, shared, or backed up simply by copying a directory. In this way, it directly supports the broader objective of democratizing access to advanced computational tools.

\item \textbf{Lilypond.} An open-source music engraving system that generates musical scores from a text-based markup language. Within this project, it has been used for the visualization of analytical results and for the generation of musical examples included throughout the research.
\end{itemize}

\section{Structure of the Rest of the Document}
The content of this project is organized into a series of chapters that follow the sequential development of the project lifecycle:

\begin{itemize}
\item \textbf{Chapter 2: State of the Art.} This chapter presents the theoretical and technological foundations that support the research. It contextualizes the so-called \textit{computational turn} within digital humanities and musicology, reviewing the principal approaches to the automated analysis of musical corpora, symbolic music representation standards such as \textit{MusicXML} and \textit{MEI}, as well as the capabilities and limitations of contemporary language models when applied to the musical domain.
\item \textbf{Chapter 3: Requirements Elicitation and Construction of the Experimental Benchmark.} This chapter describes the requirements analysis phase conducted through interviews and questionnaires with specialists in musicology and ethnomusicology. Based on this process, the real needs of researchers are identified, and a set of musicological questions is constructed to subsequently evaluate the different architectures implemented throughout the project.
\item \textbf{Chapter 4: Retrieval-Augmented Generation (RAG) Information Retrieval Systems.} This chapter examines the retrieval-augmented architectures employed as experimental baselines. It describes both the implementation based on \textit{Gemini File Search} and the local system developed using \textit{ChromaDB}, \textit{Sentence Transformers}, and \textit{Ollama}. Furthermore, it analyzes the limitations observed in both approaches when confronted with highly technical musicological queries.
\item \textbf{Chapter 5: Design and Implementation of the MCP- and SQLite-Based Hybrid Architecture.} This chapter constitutes the methodological core of the research. It presents the complete corpus-processing pipeline, including the normalization of heterogeneous musical formats, the automated extraction of metadata and musicological variables using \textit{music21}, the design of the \textit{SQLite} relational database, the implementation of the MCP server, and its integration with a local language model (\textit{Llama 3.1:8B}) for natural-language-to-SQL translation. Finally, the chapter describes the web application developed to facilitate access to the system.
\item \textbf{Chapter 6: Evaluation and Experimental Testing.} This chapter presents the validation process of the proposed solution through the systematic execution of the benchmark defined during the requirements elicitation phase. The results obtained by the different architectures are compared, analyzing response accuracy, musicological reasoning capabilities, and performance when addressing complex queries.
\item \textbf{Chapter 7: Discussion and Analysis of Limitations.} This chapter examines the results from both qualitative and quantitative perspectives. It compares the strengths and weaknesses of document-retrieval approaches against the deterministic architecture based on relational databases and MCP, analyzing phenomena such as factual hallucinations, context-window limitations, and the challenges associated with the direct interpretation of symbolic musical formats.
\item \textbf{Chapter 8: Conclusions and Future Work.} This chapter summarizes the main contributions of the research and evaluates the extent to which the proposed objectives have been achieved. It also outlines future research directions, including multi-agent architectures (\textit{Agent-to-Agent}), \textit{GraphRAG}-based systems, specialized models for musical analysis, and advanced mechanisms for exploiting semi-structured data.
\end{itemize}

The project is complemented by several appendices intended to facilitate reproducibility and support the transfer of knowledge derived from the project:

\begin{itemize}
\item \textbf{Appendix A: Knowledge Transfer and Scientific Dissemination Activities.} This appendix compiles the outputs derived from the research, including participation in national and international conferences, academic seminars, technical demonstrations, and other dissemination activities related to the project.
\item \textbf{Appendix B: Complete Benchmark and Experimental Results.} This appendix contains the full set of queries employed during the empirical evaluation, together with the responses generated by each of the four models analyzed. This material constitutes the experimental foundation of the research and enables the complete reproduction of the comparison process described in the evaluation and discussion chapters.
\end{itemize}

Finally, in order to facilitate the international dissemination of the research outcomes, the dissertation includes an English-language version of the chapters corresponding to the \textbf{Introduction} and the \textbf{Conclusions and Future Work}, which are presented respectively in Chapters 9 and 10.

\section{Availability of Source Code and Research Materials}

To promote scientific reproducibility and facilitate the dissemination of knowledge derived from this research, the source code developed throughout the project has been made publicly available in a GitHub repository.

The repository contains the complete implementation of the different architectures evaluated in this work, including the Retrieval-Augmented Generation (\textit{RAG}) systems, the infrastructure based on the \textit{Model Context Protocol} (MCP), the automated extraction pipelines for musicological variables using \textit{music21}, the database schemas employed throughout the project, and the web application developed for user interaction with the system.

The source code is available at the following repository:

\begin{center}
\url{https://github.com/marsache/MIA_TFM}
\end{center}

\section{Use of Artificial Intelligence}
This work has been developed with the occasional support of artificial intelligence tools for technical assistance and editing. Their use was limited to auxiliary tasks, such as organizing ideas, improving writing clarity, and supporting the implementation of certain parts of the code. In all cases, design decisions, result validation, and ultimate responsibility for the content and implementation rest exclusively with the author.